\documentclass[a4paper,12pt]{article}

% -------------------------------------------------
% Кодировка, язык, математика
% -------------------------------------------------
\usepackage[T2A]{fontenc}
\usepackage[utf8]{inputenc}
\usepackage[russian]{babel}
\usepackage{amsmath,amssymb,amsthm,mathtools}
\usepackage{helvet}
\renewcommand{\familydefault}{\sfdefault}
\usepackage{icomma}

% -------------------------------------------------
% Вёрстка и типографика
% -------------------------------------------------
\usepackage[
  top=18mm,
  bottom=20mm,
  left=18mm,
  right=18mm,
  headheight=15pt
]{geometry}
\usepackage{microtype}
\usepackage{setspace}
\usepackage{parskip}
\usepackage{enumitem}
\usepackage{tabularx,booktabs,longtable}
\usepackage{xcolor}
\usepackage{fancyhdr}
\usepackage{hyperref}

% -------------------------------------------------
% Графика
% -------------------------------------------------
\usepackage{tikz}
\usetikzlibrary{
  arrows.meta,
  calc,
  intersections,
  positioning
}
\usepackage{tkz-euclide}
\usepackage{pgfplots}
\pgfplotsset{compat=1.18}

% -------------------------------------------------
% Палитра
% -------------------------------------------------
\definecolor{accent}{HTML}{008080}
\definecolor{darkaccent}{HTML}{004D4D}
\definecolor{textgray}{HTML}{333333}
\definecolor{softgray}{HTML}{F1F4F4}
\definecolor{warmgray}{HTML}{777777}

% -------------------------------------------------
% Основные настройки
% -------------------------------------------------
\onehalfspacing
\setlength{\parindent}{0pt}
\setlength{\emergencystretch}{2em}
\renewcommand{\arraystretch}{1.3}

\hypersetup{
  colorlinks=true,
  linkcolor=darkaccent,
  urlcolor=darkaccent,
  pdftitle={Чек-лист: комплексное повторение ЕГЭ},
  pdfauthor={Лёвин Артём Александрович}
}

\pagestyle{fancy}
\fancyhf{}
\fancyhead[L]{\small\color{warmgray}Комплексное повторение ЕГЭ}
\fancyhead[R]{\small\color{warmgray}Николь Саркисьянц}
\fancyfoot[C]{\small\color{warmgray}\thepage}
\renewcommand{\headrulewidth}{0.4pt}
\renewcommand{\headrule}{%
  \hbox to\headwidth{\color{accent}\leaders\hrule height \headrulewidth\hfill}
}

\tkzSetUpLine[line width=1pt,color=accent]
\tkzSetUpPoint[color=accent,fill=accent,size=3]
\tkzSetUpLabel[color=black,font=\small]

% -------------------------------------------------
% Команды
% -------------------------------------------------
\DeclareMathOperator{\tg}{tg}
\DeclareMathOperator{\ctg}{ctg}
\DeclareMathOperator{\arctg}{arctg}

\newcommand{\StudentName}{Николь Саркисьянц}
\newcommand{\TeacherName}{Лёвин Артём Александрович}

\newcommand{\topic}[1]{%
  \par\bigskip
  {\Large\bfseries\color{darkaccent}#1\par}
  \vspace{0.2em}
  {\color{accent}\rule{\linewidth}{0.7pt}}
  \vspace{0.55em}
}

\newcommand{\subtopic}[1]{%
  \par\medskip
  {\large\bfseries\color{darkaccent}#1\par}
  \vspace{0.2em}
}

\newcommand{\control}[1]{%
  \begin{quote}
    \color{darkaccent}\textbf{Контрольная мысль.}\;
    \color{textgray}#1
  \end{quote}
}

\newcommand{\warning}[1]{%
  \begin{quote}
    \color{darkaccent}\textbf{Типичная ошибка.}\;
    \color{textgray}#1
  \end{quote}
}

\newlist{checklist}{itemize}{1}
\setlist[checklist]{
  label=\(\square\),
  leftmargin=1.7em,
  itemsep=0.35em,
  topsep=0.35em
}

\setlist[enumerate]{
  leftmargin=2em,
  itemsep=0.45em,
  topsep=0.4em
}

% =================================================
\begin{document}

% =================================================
% Титульный лист
% =================================================
\begin{titlepage}
  \thispagestyle{empty}
  \centering

  \vspace*{25mm}

  {\color{accent}\rule{0.42\linewidth}{1.2pt}\par}
  \vspace{10mm}

  {\fontsize{25}{30}\selectfont
    \bfseries\color{darkaccent}
    Чек-лист
  \par}

  \vspace{5mm}

  {\fontsize{19}{24}\selectfont
    \bfseries
    Комплексное повторение задач ЕГЭ
  \par}

  \vspace{5mm}

  {\large
    Геометрия, векторы, вероятность, степени, логарифмы,\\
    задачи на работу, показательная функция и тригонометрия
  \par}

  \vspace{18mm}

  {\color{accent}\rule{0.18\linewidth}{0.8pt}\par}

  \vspace{14mm}

  {\large
    \TeacherName
  \par}

  \vspace{3mm}

  {\large
    эксклюзивно для
  \par}

  \vspace{3mm}

  {\LARGE\bfseries\color{darkaccent}
    \StudentName
  \par}

  \vfill

  {\large ЕГЭ по профильной математике\par}
  \vspace{4mm}
  {\large \today\par}

  \begin{tikzpicture}[remember picture,overlay]
    \draw[
      darkaccent!45,
      line width=1.1pt
    ]
    ([xshift=12mm,yshift=18mm]current page.south west)
    .. controls
    ([xshift=55mm,yshift=31mm]current page.south west)
    and
    ([xshift=-62mm,yshift=9mm]current page.south east)
    ..
    ([xshift=-12mm,yshift=20mm]current page.south east);

    \draw[
      accent!35,
      line width=0.6pt
    ]
    ([xshift=18mm,yshift=13mm]current page.south west)
    .. controls
    ([xshift=72mm,yshift=25mm]current page.south west)
    and
    ([xshift=-78mm,yshift=6mm]current page.south east)
    ..
    ([xshift=-18mm,yshift=15mm]current page.south east);
  \end{tikzpicture}
\end{titlepage}

% =================================================
% Основной чек-лист
% =================================================
\topic{Карта повторения}

\begin{table}[htbp]
  \caption{Основные задачи занятия}
  \centering
  \begin{tabularx}{\linewidth}{@{}p{0.25\linewidth}X@{}}
    \toprule
    \textbf{Раздел} & \textbf{Что требуется уметь} \\
    \midrule
    Планиметрия &
    Использовать свойства равнобедренного треугольника и теорему о внешнем угле. \\

    Векторы &
    Находить координаты вектора, выполнять линейные операции и вычислять длину. \\

    Вероятность &
    Строить пространство исходов, работать с противоположными событиями, объединением и пересечением. \\

    Степени и корни &
    Представлять числа в степенном виде и применять свойства степеней. \\

    Логарифмы &
    Проверять условия существования и выбирать подходящее логарифмическое тождество. \\

    Задачи на работу &
    Составлять таблицу «производительность --- время --- работа» и получать уравнение. \\

    Показательная функция &
    Определять основание функции по точке графика и вычислять значение функции. \\

    Тригонометрические уравнения &
    Сводить показательное уравнение к тригонометрическому и сохранять все серии корней. \\
    \bottomrule
  \end{tabularx}
\end{table}

\begin{checklist}
  \item Я записываю обозначения до начала вычислений.
  \item Я проверяю знак каждой координаты, степени и дроби.
  \item Я выбираю формулу по структуре выражения.
  \item Я сохраняю корни при разложении выражения на множители.
  \item Я проверяю ОДЗ исходного выражения.
  \item Я оцениваю правдоподобие ответа.
\end{checklist}

% -------------------------------------------------
\topic{1. Внешний угол равнобедренного треугольника}

Пусть в равнобедренном треугольнике \(ABC\) основанием служит
\(AB\). Тогда боковые стороны равны:
\[
AC=BC,
\]
а углы при основании равны:
\[
\angle BAC=\angle ABC.
\]

Внешний угол при вершине \(C\) равен сумме двух внутренних углов,
которые с ним не смежны:
\[
\angle BCF=\angle BAC+\angle ABC.
\]

\begin{figure}[htbp]
  \centering
  \begin{tikzpicture}[scale=1]
    \tkzDefPoint(-5.05,0){A}
    \tkzDefPoint(5.05,0){B}
    \tkzDefPoint(0,0.8){C}
    \tkzDefPoint(2.525,1.2){F}

    \tkzDrawSegments(A,B B,C C,A C,F)
    \tkzDrawPoints(A,B,C,F)

    \tkzMarkSegments[mark=|](A,C B,C)
    \tkzMarkAngle[size=0.8](B,C,F)
    \tkzLabelAngle[pos=1.25](B,C,F){$18^\circ$}

    \tkzMarkAngle[size=0.8](C,A,B)
    \tkzMarkAngle[size=0.8](A,B,C)
    \tkzLabelAngle[pos=1.25](C,A,B){$x$}
    \tkzLabelAngle[pos=1.25](A,B,C){$x$}

    \tkzLabelPoints[below](A,B)
    \tkzLabelPoints[above left](C)
    \tkzLabelPoints[above](F)
  \end{tikzpicture}
  \caption{Внешний угол при вершине равен сумме углов при основании}
\end{figure}

Если
\[
\angle BCF=18^\circ,
\]
то
\[
2x=18^\circ,
\qquad
x=9^\circ.
\]

Возможен и полный расчёт через внутренний угол:
\[
\angle ACB=180^\circ-18^\circ=162^\circ,
\]
\[
\angle ABC
=
\frac{180^\circ-162^\circ}{2}
=
9^\circ.
\]

\control{Короткое решение через теорему о внешнем угле экономит время. Полный расчёт удобно использовать для самопроверки.}

\begin{checklist}
  \item Я определил основание равнобедренного треугольника.
  \item Я отметил равные углы при основании.
  \item Я отличаю внутренний угол от смежного внешнего угла.
  \item Сумма внутренних углов моего треугольника равна \(180^\circ\).
\end{checklist}

% -------------------------------------------------
\topic{2. Координаты и длина вектора}

Если вектор смещает точку на \(x\) единиц по горизонтали и на
\(y\) единиц по вертикали, его координаты записываются так:
\[
\vec a=(x_a;y_a).
\]

Направление учитывается знаком:

\[
\begin{array}{c|c}
\text{направление} & \text{знак координаты} \\ \hline
\text{вправо или вверх} & + \\
\text{влево или вниз} & -
\end{array}
\]

Для векторов
\[
\vec a=(x_a;y_a),
\qquad
\vec b=(x_b;y_b)
\]
выполняются правила:
\[
\vec a-\vec b
=
(x_a-x_b;\,y_a-y_b),
\]
\[
k\vec a=(kx_a;\,ky_a).
\]

Длина вектора \(\vec c=(x_c;y_c)\):
\[
|\vec c|
=
\sqrt{x_c^2+y_c^2}.
\]

\subtopic{Пример}

Пусть
\[
\vec a=(7;0),
\qquad
\vec b=(-5;-5).
\]

Обозначим
\[
\vec c=\vec a-\vec b.
\]

Тогда
\[
\vec c
=
\bigl(7-(-5);\,0-(-5)\bigr)
=
(12;5),
\]
\[
|\vec c|
=
\sqrt{12^2+5^2}
=
\sqrt{169}
=
13.
\]

\begin{figure}[htbp]
  \centering
  \begin{tikzpicture}[x=0.42cm,y=0.42cm,>=Stealth]
    \draw[step=1,softgray] (-6,-6) grid (13,6);
    \draw[->,line width=0.8pt] (-6,0)--(13.5,0)
      node[right] {$x$};
    \draw[->,line width=0.8pt] (0,-6)--(0,6.5)
      node[above] {$y$};

    \draw[->,very thick,accent]
      (0,0)--(7,0)
      node[midway,above] {$\vec a=(7;0)$};

    \draw[->,very thick,darkaccent]
      (0,0)--(-5,-5)
      node[midway,above left] {$\vec b=(-5;-5)$};

    \draw[->,very thick,textgray]
      (0,0)--(12,5)
      node[midway,above] {$\vec a-\vec b=(12;5)$};

    \fill[black] (0,0) circle (2pt);
  \end{tikzpicture}
  \caption{Координаты вектора определяются его направленным смещением}
\end{figure}

\warning{Запись \(2\vec a\) означает умножение каждой координаты вектора на \(2\). Запись \(2|\vec a|\) означает удвоенную длину вектора.}

\begin{checklist}
  \item Я учитываю направление смещения по каждой оси.
  \item При вычитании векторов я вычитаю соответствующие координаты.
  \item При умножении вектора на число я умножаю обе координаты.
  \item Под корнем в формуле длины стоит сумма квадратов координат.
\end{checklist}

% -------------------------------------------------
\topic{3. Теория вероятностей}

\subtopic{3.1. Таблица равновозможных исходов}

Перед каждым из двух матчей команда с вероятностью
\(\frac12\) начинает игру с мячом.

Обозначим:
\[
+\; \text{--- команда начинает игру с мячом},
\qquad
-\; \text{--- команда не начинает игру с мячом}.
\]

Пространство исходов:
\[
(++),\quad (+-),\quad (-+),\quad (--).
\]

Все четыре исхода равновероятны. Событию «команда ни разу не начнёт
игру с мячом» соответствует один исход:
\[
(--).
\]

Следовательно,
\[
P=\frac14=0,25.
\]

\control{Обозначения должны отражать событие из условия. Это снижает риск неверной интерпретации орла и решки.}

\subtopic{3.2. Противоположное событие}

Для события \(A\) противоположное событие обозначается
\(\overline A\).

Всегда выполняется:
\[
P(\overline A)=1-P(A).
\]

Пример:
\[
P(A)=0,2
\quad\Longrightarrow\quad
P(\overline A)=1-0,2=0,8.
\]

\subtopic{3.3. Объединение и пересечение событий}

Обозначения:
\[
A\cap B
\quad\text{--- произошли оба события},
\]
\[
A\cup B
\quad\text{--- произошло хотя бы одно из событий}.
\]

Формула сложения вероятностей:
\[
P(A\cup B)
=
P(A)+P(B)-P(A\cap B).
\]

Пересечение вычитается один раз, поскольку при сложении
\(P(A)+P(B)\) оно учитывается дважды.

\subtopic{Пример с кофейными автоматами}

Пусть:
\[
A=\{\text{кофе закончился в первом автомате}\},
\]
\[
B=\{\text{кофе закончился во втором автомате}\}.
\]

По условию:
\[
P(A)=P(B)=0,2,
\qquad
P(A\cap B)=0,16.
\]

Требуется найти вероятность того, что кофе останется в обоих
автоматах:
\[
P(\overline A\cap\overline B).
\]

Событие \(\overline A\cap\overline B\) противоположно событию
\(A\cup B\). Поэтому
\[
P(\overline A\cap\overline B)
=
1-P(A\cup B).
\]

Вычислим:
\[
\begin{aligned}
P(A\cup B)
&=
0,2+0,2-0,16\\
&=
0,24,
\end{aligned}
\]
\[
P(\overline A\cap\overline B)
=
1-0,24
=
0,76.
\]

\begin{figure}[htbp]
  \centering
  \begin{tikzpicture}[scale=1.05]
    \fill[softgray] (0,0) rectangle (7,4);
    \draw[line width=0.9pt] (0,0) rectangle (7,4);

    \fill[accent!10] (2.7,2) circle (1.35);
    \fill[darkaccent!10] (4.3,2) circle (1.35);

    \begin{scope}
      \clip (2.7,2) circle (1.35);
      \fill[accent!30] (4.3,2) circle (1.35);
    \end{scope}

    \draw[accent,line width=1pt] (2.7,2) circle (1.35);
    \draw[darkaccent,line width=1pt] (4.3,2) circle (1.35);

    \node[accent] at (2.05,3.1) {$A$};
    \node[darkaccent] at (4.95,3.1) {$B$};
    \node at (3.5,2) {$0,16$};
    \node at (1.8,2) {$0,04$};
    \node at (5.2,2) {$0,04$};
    \node at (3.5,0.55) {$0,76$};
    \node[anchor=north east] at (6.9,3.9) {$\Omega$};
  \end{tikzpicture}
  \caption{Распределение вероятностей между областями диаграммы Эйлера}
\end{figure}

\warning{Круг \(A\) должен соответствовать вероятности события \(A\). В данном примере площадь каждого круга равна \(0,2\), поскольку круг означает окончание кофе.}

\begin{checklist}
  \item Я ввёл события буквами и расшифровал их.
  \item Я различаю \(A\cap B\) и \(A\cup B\).
  \item Я вычитаю пересечение в формуле объединения.
  \item Итоговая вероятность принадлежит отрезку \([0;1]\).
\end{checklist}

% -------------------------------------------------
\topic{4. Степени и корни}

Основные свойства степеней:
\[
a^m\cdot a^n=a^{m+n},
\]
\[
\frac{a^m}{a^n}=a^{m-n},
\qquad a\ne0,
\]
\[
\left(a^m\right)^n=a^{mn},
\]
\[
a^{-n}=\frac1{a^n},
\qquad a\ne0,
\]
\[
\sqrt[n]{a}=a^{1/n},
\]
если выражение определено.

Для нечётного \(n\) корень \(\sqrt[n]{a}\) существует при любом
действительном \(a\).

\subtopic{Пример}

Решим уравнение:
\[
x-13=\sqrt[3]{64}.
\]

Представим число \(64\) в степенном виде:
\[
64=2^6.
\]

Тогда
\[
\sqrt[3]{64}
=
64^{1/3}
=
\left(2^6\right)^{1/3}
=
2^2
=
4.
\]

Следовательно,
\[
x-13=4,
\qquad
x=17.
\]

\control{Для работы с корнями и степенями составное число удобно раскладывать на простые множители.}

\begin{checklist}
  \item Я различаю показатель степени и основание.
  \item При возведении степени в степень я перемножаю показатели.
  \item Отрицательный показатель превращает выражение в обратное.
  \item Для корня чётной степени я проверяю неотрицательность подкоренного выражения.
\end{checklist}

% -------------------------------------------------
\topic{5. Логарифмы}

\subtopic{5.1. Условия существования}

Выражение
\[
\log_a b
\]
определено при условиях:
\[
a>0,
\qquad
a\ne1,
\qquad
b>0.
\]

\subtopic{5.2. Основные формулы}

При выполнении условий существования:
\[
\log_a(bc)
=
\log_a b+\log_a c,
\]
\[
\log_a\frac bc
=
\log_a b-\log_a c,
\]
\[
\log_a(b^n)
=
n\log_a b,
\]
\[
\log_{a^n}b
=
\frac1n\log_a b,
\qquad n\ne0,
\]
\[
a^{\log_a b}=b,
\]
\[
\log_a b
=
\frac1{\log_b a},
\]
\[
\log_a b
=
\frac{\log_c b}{\log_c a},
\qquad
c>0,\quad c\ne1.
\]

\subtopic{5.3. Пример на основное логарифмическое тождество}

Вычислим:
\[
7^{2+\log_7 2}.
\]

Используем свойство
\[
a^{m+n}=a^m\cdot a^n:
\]
\[
7^{2+\log_7 2}
=
7^2\cdot 7^{\log_7 2}.
\]

По основному логарифмическому тождеству:
\[
7^{\log_7 2}=2.
\]

Следовательно,
\[
7^{2+\log_7 2}
=
49\cdot2
=
98.
\]

\warning{Формула \(a^{\log_a b}=b\) применяется при совпадении основания степени и основания логарифма. Перед логарифмом в показателе должен стоять коэффициент \(1\).}

\begin{checklist}
  \item Я проверил основание и аргумент каждого логарифма.
  \item При сложении логарифмов аргументы перемножаются.
  \item При вычитании логарифмов аргументы делятся.
  \item Степень аргумента становится множителем перед логарифмом.
  \item Я распознаю конструкцию \(a^{\log_a b}\).
\end{checklist}

% -------------------------------------------------
\topic{6. Задачи на совместную работу}

Для задач на работу используется связь:
\[
\text{работа}
=
\text{производительность}\cdot\text{время}.
\]

Если вся работа принимается за \(1\), а исполнитель выполняет её
за \(t\) минут, его производительность равна:
\[
\frac1t.
\]

\subtopic{Пример}

Катя и Настя вместе пропалывают грядку за \(15\) минут.
Настя выполняет эту работу за \(40\) минут. Требуется найти время
работы Кати.

Обозначим время работы Кати через \(x\) минут.

\begin{table}[htbp]
  \caption{Модель задачи на работу}
  \centering
  \begin{tabularx}{\linewidth}{@{}Xccc@{}}
    \toprule
    \textbf{Исполнитель} &
    \textbf{Производительность} &
    \textbf{Время} &
    \textbf{Работа} \\
    \midrule
    Настя & \(\dfrac1{40}\) & \(40\) & \(1\) \\
    Катя & \(\dfrac1x\) & \(x\) & \(1\) \\
    Катя и Настя & \(\dfrac1{15}\) & \(15\) & \(1\) \\
    \bottomrule
  \end{tabularx}
\end{table}

При совместной работе производительности складываются:
\[
\frac1{40}+\frac1x=\frac1{15}.
\]

Приведём дроби к общему знаменателю:
\[
\frac1x
=
\frac1{15}-\frac1{40}
=
\frac8{120}-\frac3{120}
=
\frac5{120}
=
\frac1{24}.
\]

Следовательно,
\[
x=24.
\]

Катя выполняет работу за \(24\) минуты.

\control{Рабочее уравнение составляется по столбцу производительности.}

\begin{checklist}
  \item Я принял всю работу за \(1\).
  \item Я выразил производительность через время.
  \item При совместной работе я сложил производительности.
  \item Я проверил, что совместное время меньше времени каждого исполнителя.
\end{checklist}

% -------------------------------------------------
\topic{7. Показательная функция по графику}

Показательная функция имеет вид:
\[
f(x)=a^x,
\]
где
\[
a>0,
\qquad
a\ne1.
\]

Любой график функции \(f(x)=a^x\) проходит через точку
\[
(0;1),
\]
поскольку
\[
a^0=1.
\]

Эта точка подтверждает общий вид функции, однако значение \(a\)
по ней определить невозможно.

Для нахождения основания выбирают точку с ненулевой координатой
\(x\).

\subtopic{Пример}

График проходит через точку
\[
B(-1;3).
\]

Подставим координаты:
\[
3=a^{-1}.
\]

Поскольку
\[
a^{-1}=\frac1a,
\]
получаем:
\[
\frac1a=3,
\qquad
a=\frac13.
\]

Следовательно,
\[
f(x)=\left(\frac13\right)^x.
\]

Найдём:
\[
f(-3)
=
\left(\frac13\right)^{-3}
=
3^3
=
27.
\]

\begin{figure}[htbp]
  \centering
  \begin{tikzpicture}
    \begin{axis}[
      width=9cm,
      height=7cm,
      axis lines=middle,
      unit vector ratio=1 1,
      xmin=-1.4,
      xmax=2.6,
      ymin=-0.2,
      ymax=4.2,
      samples=250,
      clip=false,
      xlabel={$x$},
      ylabel={$y$},
      xtick={-1,0,1,2},
      ytick={1,2,3,4},
      grid=both,
      minor tick num=1
    ]
      \addplot[
        domain=-1.3:2.5,
        very thick,
        color=accent
      ]
      {exp(ln(1/3)*x)};

      \addplot[
        only marks,
        mark=*,
        mark size=2.5pt,
        color=darkaccent
      ]
      coordinates {(-1,3) (0,1)};

      \node[anchor=south west] at (axis cs:-1,3)
        {$B(-1;3)$};
      \node[anchor=south west] at (axis cs:0,1)
        {$(0;1)$};
    \end{axis}
  \end{tikzpicture}
  \caption{График функции \(f(x)=\left(\frac13\right)^x\)}
\end{figure}

\begin{checklist}
  \item Я выбрал узловую точку с ненулевой координатой \(x\).
  \item Я подставил координаты точки в равенство \(y=a^x\).
  \item Я учёл правило отрицательной степени.
  \item Найденное основание удовлетворяет условиям \(a>0\) и \(a\ne1\).
\end{checklist}

% -------------------------------------------------
\topic{8. Показательно-тригонометрическое уравнение}

Рассмотрим уравнение:
\[
7^{2\cos x}=49^{\sin 2x}.
\]

\subtopic{Шаг 1. Приведение оснований}

Поскольку
\[
49=7^2,
\]
получаем:
\[
7^{2\cos x}
=
\left(7^2\right)^{\sin 2x}
=
7^{2\sin 2x}.
\]

Основания равны и принадлежат промежутку
\[
7>0,\qquad 7\ne1.
\]

Функция \(7^t\) строго возрастает, поэтому показатели равны:
\[
2\cos x=2\sin 2x.
\]

\subtopic{Шаг 2. Приведение углов}

Перенесём всё в левую часть:
\[
2\cos x-2\sin 2x=0.
\]

Используем формулу двойного угла:
\[
\sin 2x=2\sin x\cos x.
\]

Тогда:
\[
2\cos x-4\sin x\cos x=0.
\]

Вынесем общий множитель:
\[
2\cos x(1-2\sin x)=0.
\]

\subtopic{Шаг 3. Совокупность уравнений}

Произведение равно нулю, если хотя бы один множитель равен нулю:
\[
\left[
\begin{aligned}
\cos x&=0,\\
1-2\sin x&=0.
\end{aligned}
\right.
\]

Первое уравнение:
\[
\cos x=0
\quad\Longrightarrow\quad
x=\frac\pi2+\pi n,
\qquad n\in\mathbb Z.
\]

Второе уравнение:
\[
\sin x=\frac12.
\]

Отсюда:
\[
x=\frac\pi6+2\pi n,
\qquad n\in\mathbb Z,
\]
или
\[
x=\frac{5\pi}6+2\pi n,
\qquad n\in\mathbb Z.
\]

Итог:
\[
\boxed{
x=\frac\pi2+\pi n
\quad\text{или}\quad
x=\frac\pi6+2\pi n
\quad\text{или}\quad
x=\frac{5\pi}6+2\pi n,
\quad n\in\mathbb Z
}
\]

Исходное уравнение определено при всех
\[
x\in\mathbb R.
\]

\warning{Деление уравнения на \(\cos x\) удалило бы серию корней \(x=\frac\pi2+\pi n\). Безопасный приём состоит в вынесении общего множителя и переходе к совокупности.}

\begin{checklist}
  \item Я привёл основания степеней к одному числу.
  \item Я корректно перемножил показатели в конструкции \((a^m)^n\).
  \item Я привёл тригонометрические функции к одинаковому углу.
  \item Я вынес общий множитель без сокращения на выражение с переменной.
  \item Я записал все серии корней и указал \(n\in\mathbb Z\).
\end{checklist}

% -------------------------------------------------
\topic{9. Полезные формулы двойного угла}

\[
\sin 2\alpha
=
2\sin\alpha\cos\alpha.
\]

\[
\cos 2\alpha
=
\cos^2\alpha-\sin^2\alpha.
\]

\[
\cos 2\alpha
=
2\cos^2\alpha-1.
\]

\[
\cos 2\alpha
=
1-2\sin^2\alpha.
\]

Для уравнения
\[
\sin x=\frac12
\]
табличное значение:
\[
\arcsin\frac12=\frac\pi6.
\]

Две серии решений:
\[
x=\frac\pi6+2\pi n,
\]
\[
x=\pi-\frac\pi6+2\pi n
=
\frac{5\pi}6+2\pi n,
\qquad n\in\mathbb Z.
\]

\subtopic{Табличные значения синуса}

\begin{table}[htbp]
  \caption{Значения синуса основных углов}
  \centering
  \begin{tabularx}{\linewidth}{@{}Xccccc@{}}
    \toprule
    \(x\) &
    \(0\) &
    \(\dfrac\pi6\) &
    \(\dfrac\pi4\) &
    \(\dfrac\pi3\) &
    \(\dfrac\pi2\) \\
    \midrule
    \(\sin x\) &
    \(0\) &
    \(\dfrac12\) &
    \(\dfrac{\sqrt2}{2}\) &
    \(\dfrac{\sqrt3}{2}\) &
    \(1\) \\
    \bottomrule
  \end{tabularx}
\end{table}

% -------------------------------------------------
\topic{10. Процентный множитель в финансовых задачах}

При увеличении величины на \(p\%\) используется множитель:
\[
q=1+\frac{p}{100}.
\]

Например:
\[
p=8\%
\quad\Longrightarrow\quad
q=1+\frac8{100}=1,08.
\]

При уменьшении величины на \(p\%\):
\[
q=1-\frac{p}{100}.
\]

В задачах о кредитах удобно:

\begin{enumerate}
  \item ввести переменную \(x\) для платежа;
  \item построить таблицу по периодам;
  \item последовательно учитывать начисление процентов и платёж;
  \item получить формулу, связывающую исходную сумму, ставку и платежи;
  \item проверить единицы измерения и процентный множитель.
\end{enumerate}

\warning{Запись \(1,8\) соответствует увеличению на \(80\%\). Для ставки \(8\%\) используется множитель \(1,08\).}

% -------------------------------------------------
\topic{Итоговый алгоритм самопроверки}

\begin{enumerate}
  \item Перепишите условие в математических обозначениях.
  \item Выделите известные и искомые величины.
  \item Запишите условия существования выражений.
  \item Выберите основную формулу или теорему.
  \item Выполняйте преобразования по одному логическому шагу.
  \item Следите за знаками, скобками и порядком вычитания.
  \item После разложения на множители решите каждое уравнение совокупности.
  \item Подставьте ответ в исходную модель или оцените его смысл.
  \item Проверьте единицы измерения.
  \item Запишите ответ в требуемом формате.
\end{enumerate}

% =================================================
% Тренировочный блок
% =================================================
\clearpage
\thispagestyle{fancy}

\topic{Тренировочный блок}

Во всех заданиях записывайте основные формулы и ключевые этапы
решения. Ответы расположены на следующей странице.

\subtopic{Средний уровень}

\begin{enumerate}
  \item
  В равнобедренном треугольнике \(ABC\) основанием является
  \(AB\). Внешний угол при вершине \(C\) равен \(26^\circ\).
  Найдите \(\angle ABC\).

  \item
  Даны векторы
  \[
  \vec a=(4;-3),
  \qquad
  \vec b=(-2;1).
  \]
  Найдите длину вектора
  \[
  \vec a-\vec b.
  \]

  \item
  Перед каждым из двух матчей команда с вероятностью
  \(\frac12\) начинает игру с мячом. Найдите вероятность того,
  что команда ни в одном матче не начнёт игру с мячом.
\end{enumerate}

\subtopic{Уровень выше среднего}

\begin{enumerate}[resume]
  \item
  В торговом центре установлены два кофейных автомата.
  Вероятность окончания кофе в каждом автомате равна \(0,25\).
  Вероятность окончания кофе сразу в обоих автоматах равна
  \(0,10\). Найдите вероятность того, что кофе останется в обоих
  автоматах.

  \item
  Решите уравнение:
  \[
  x-9=\sqrt[3]{125}.
  \]

  \item
  Вычислите:
  \[
  5^{2+\log_5 4}.
  \]

  \item
  Первый исполнитель выполняет работу за \(20\) минут.
  Вместе со вторым исполнителем они выполняют эту работу за
  \(12\) минут. За сколько минут второй исполнитель выполнит
  работу самостоятельно?

  \item
  Показательная функция
  \[
  f(x)=a^x
  \]
  проходит через точку
  \[
  (-2;9).
  \]
  Найдите
  \[
  f(-3).
  \]

  \item
  Даны векторы
  \[
  \vec a=(2;-1),
  \qquad
  \vec b=(-3;4).
  \]
  Найдите длину вектора
  \[
  2\vec a-\vec b.
  \]
\end{enumerate}

\subtopic{Сложный уровень}

\begin{enumerate}[resume]
  \item
  Решите уравнение
  \[
  3^{2\cos x}=9^{\sin 2x}
  \]
  на промежутке
  \[
  x\in[0;2\pi).
  \]
\end{enumerate}

\vfill

\begin{checklist}
  \item В геометрической задаче я использовал равенство углов при основании.
  \item В векторных задачах я выполнил операции по координатам.
  \item В вероятностной задаче я проверил принадлежность ответа отрезку \([0;1]\).
  \item В задаче на работу я сравнил совместное и индивидуальное время.
  \item В показательном уравнении я сохранил все множители.
\end{checklist}

% =================================================
% Ответы
% =================================================
\clearpage
\thispagestyle{fancy}

\topic{Ответы к тренировочному блоку}

\begin{table}[htbp]
  \caption{Краткие ответы и контрольные идеи}
  \centering
  \begin{tabularx}{\linewidth}{@{}p{0.07\linewidth}p{0.24\linewidth}X@{}}
    \toprule
    \textbf{\(\№\)} &
    \textbf{Ответ} &
    \textbf{Контрольная идея} \\
    \midrule

    1 &
    \(13^\circ\) &
    Внешний угол равен сумме двух равных углов при основании. \\

    2 &
    \(2\sqrt{13}\) &
    \(\vec a-\vec b=(6;-4)\), затем применяется формула длины. \\

    3 &
    \(0,25\) &
    Подходит один исход из четырёх равновероятных исходов. \\

    4 &
    \(0,60\) &
    \(1-(0,25+0,25-0,10)\). \\

    5 &
    \(14\) &
    \(\sqrt[3]{125}=5\). \\

    6 &
    \(100\) &
    \(5^2\cdot5^{\log_5 4}=25\cdot4\). \\

    7 &
    \(30\) минут &
    \(\dfrac1{20}+\dfrac1x=\dfrac1{12}\). \\

    8 &
    \(27\) &
    \(a^{-2}=9\), поэтому \(a=\dfrac13\). \\

    9 &
    \(\sqrt{85}\) &
    \(2\vec a-\vec b=(7;-6)\). \\

    10 &
    \(\dfrac\pi6,\;
      \dfrac\pi2,\;
      \dfrac{5\pi}6,\;
      \dfrac{3\pi}2\) &
    После приведения оснований:
    \(2\cos x(1-2\sin x)=0\). \\

    \bottomrule
  \end{tabularx}
\end{table}

\vspace{8mm}

\begin{center}
  {\large\bfseries\color{darkaccent}
    Основной ориентир: точность обозначений,
    сохранение всех корней и обязательная проверка ответа.
  \par}

  \vspace{6mm}

  {\color{accent}\rule{0.55\linewidth}{0.8pt}}
\end{center}

\end{document}